\chapter{Deferred foundations (roadmap)}

This chapter makes explicit the infrastructure the rest of the blueprint rests on but
does not build. Each node below is a \emph{foundation}: a body of mathematics the paper
cites as known, or a Mathlib gap, encoded in the Lean development as a \texttt{class} of
assumed properties with no instance. The scope-(a) content of Sections~2, 3 and~5 is
stated faithfully over these assumptions; the dependency edges from those content nodes
into this chapter are what turn the dependency graph into a roadmap of the unbuilt walls.

\emph{Reading the colours.}
\begin{itemize}
  \item \textbf{Green} (dark green fill) is the finish line: statement and proof are
    formalized \emph{and every node it depends on, transitively, is too}, so it rests on
    nothing assumed. Green means full formalization including foundations, relative to
    nothing assumed.
  \item \textbf{Scaffold} (green border, possibly the lighter green or blue fill) is a
    faithfully-encoded statement, perhaps with a compiling proof, that still depends on at
    least one orange foundation. A compiling proof keeps its green; the ``not yet on solid
    ground'' signal is the dependency edge into an orange node, not a withheld colour.
  \item \textbf{Orange} (\notready) is a buildable foundation not yet formalized, or a
    content node whose statement itself is a \texttt{True}-placeholder. These are the
    grind queue.
  \item \textbf{Grey, dashed} is a \emph{deliberately abstract} foundation: a parametric
    assumption that is never instantiated by design (see below), not a TODO.
\end{itemize}
The project is finished when every node is \textbf{green, or a foundation deliberately
left abstract with a recorded reason} --- not necessarily 100\% green.

\emph{Tooling caveat.} The dark-green ``fully formalized'' colour treats unbuilt
\emph{definition}-foundations as transparent to its computation, so a proved result
resting only on them can mis-render dark green; whenever a node newly turns dark green,
re-check that its foundation ancestry is genuinely built before trusting the colour.

\section{Deliberately-abstract parametric wrappers}

A field theory is a functor out of an \emph{arbitrary} cobordism geometry, in the same way
an Atiyah-style TQFT is a functor on an arbitrary cobordism category. The
arbitrary-geometry parametrization, and the tangent / dual / globally-hyperbolic plumbing
threaded through it, are therefore not objects to be constructed: staying parametric is the
intended permanent state. They are recorded here, grey and dashed, so the buildable-orange
queue reads unambiguously.

\begin{definition}[Arbitrary-cobordism-geometry parametrization]
  \label{abstract:cobordism-domain}
  \abstractnode
  The TQFT-style abstraction underlying Section~3: a field theory is developed over an
  arbitrary \texttt{CobordismGeometry}, never a concrete one. This node carries the
  deliberately-abstract aspect of \texttt{CobordismGeometry} (its bare object/morphism
  parametrization, as opposed to the smooth geometry of
  Definition~\ref{found:smooth-cobordism-geometry}), together with the abstract
  \texttt{TangentStructure} (an abstract manifold-like type with a tangent space) and the
  \texttt{DualConjugateGeometry} operations $\Sigma \mapsto \Sigma^*$, $\Sigma \mapsto
  \bar\Sigma$. These are parametric plumbing: a concrete instance is never a goal, so the
  node is abstract, not orange.
\end{definition}

\begin{definition}[Arbitrary-Lorentzian-geometry parametrization]
  \label{abstract:lorentzian-domain}
  \abstractnode
  The Section~5 counterpart: the development is parametric over an arbitrary
  \texttt{LorentzianCobordismGeometry}, never a concrete one. This node carries that
  abstraction together with the \texttt{GhCategory} closure/openness assumptions (that the
  globally hyperbolic cobordisms are closed under concatenation and form an open
  subcategory), which are categorical bookkeeping over the abstract geometry rather than a
  construction to formalize.
\end{definition}

\section{Buildable foundations}

These are the walls to build. The existing nodes
Definition~\ref{def:nuclear-frechet} (nuclear Fr\'echet spaces),
Definition~\ref{def:metc-complex-manifold} ($\mathrm{Met}_{\mathbb C}(M)$ as a complex
manifold) and Definition~\ref{def:holomorphic-bundle} (holomorphic vector bundles) are
themselves foundations of this kind, already orange in Sections~3; they are not repeated
here. Difficulty is flagged as \emph{grind} (sit down and formalize),
\emph{Mathlib-scale} (a library to contribute) or \emph{research-scale} (a genuine wall).

\begin{definition}[Limits and completed tensor products of nuclear spaces]
  \label{found:tvs-limits-tensor}
  \lean{CountableNuclearInverseSystem, HasNuclearInverseLimit,
    CountableNuclearDirectSystem, HasNuclearDirectLimit,
    HasCompletedNuclearTensor}
  \notready
  \uses{def:nuclear-frechet}
  Countable inverse and direct limits of nuclear maps of Fr\'echet spaces, and the unique
  completed tensor product of nuclear spaces. This is what actually realizes
  $\hat E_\Sigma = \varprojlim E_{\Sigma''}$, $\check E_\Sigma = \varinjlim E_{\Sigma'}$ and
  $E \otimes F$, all of which Section~3 asserts but defers (the ``limit construction is
  deferred'' in \texttt{FieldTheory.lean}; the tensor in \texttt{TensorAxiom.lean}). It is
  the second Mathlib gap named in the README. Missing: limits and completed tensor products
  of topological vector spaces (Mathlib has only algebraic tensor products and finite-rank
  traces). \emph{Mathlib-scale.}
\end{definition}

\begin{definition}[The smooth cobordism geometry]
  \label{found:smooth-cobordism-geometry}
  \notready
  \uses{def:allowable}
  The geometric content abstracted by \texttt{CobordismGeometry}: germs of $d$-manifolds
  along closed co-oriented $(d-1)$-submanifolds, complex metrics as sections of
  $\mathrm{Met}_{\mathbb C}$, cobordisms up to boundary-fixing isometry, their gluing, and
  the disjoint unions of \texttt{MonoidalCobordism}. The categorical skeleton over this
  geometry (the \texttt{Semicategory}) is already built; what is missing is the
  smooth-manifold geometry itself, none of which (manifold germs, gluing) is in Mathlib.
  This geometry conceptually rests on the complex-manifold structure
  $\mathrm{Met}_{\mathbb C}$ (Definition~\ref{def:metc-complex-manifold}), of which its
  metrics are sections; the dependency edge to that node is deliberately omitted to keep
  the dependency graph acyclic, since Definition~\ref{def:metc-complex-manifold} already
  depends on Definition~\ref{def:cobordism-category} and the edge would close a cycle. The
  recorded dependency on Definition~\ref{def:allowable} instead reflects that those
  cobordism metrics are allowable complex metrics.
  Pairs with the deliberately-abstract Definition~\ref{abstract:cobordism-domain}.
  \emph{Research-scale.}
\end{definition}

\begin{definition}[Lorentzian and causal geometry]
  \label{found:lorentzian-causal-geometry}
  \notready
  \uses{found:smooth-cobordism-geometry}
  The geometric content abstracted by \texttt{LorentzianCobordismGeometry},
  \texttt{LightConeGeometry} and \texttt{GhCategory}: Lorentzian metrics of signature
  $(d-1,1)$, light-cone structure, timelike geodesics, smooth time-functions, global
  hyperbolicity, and the geodesic normal form $h_t - \mathrm dt^2$. Lorentzian causal
  geometry is standard mathematics but absent from Mathlib. Pairs with the
  deliberately-abstract Definition~\ref{abstract:lorentzian-domain}.
  \emph{Research-scale.}
\end{definition}

\begin{definition}[Real-analytic structure and holomorphic complexification]
  \label{found:real-analytic-complexification}
  \notready
  \uses{found:lorentzian-causal-geometry, def:metc-complex-manifold}
  The real-analytic structure and its holomorphic complexification: a real-analytic
  manifold $M$ as a totally real submanifold of a complex manifold $M_{\mathbb C}$, with a
  holomorphic symmetric form $g$ on $TM_{\mathbb C}$ inducing the allowable metric, and the
  isotopies of $M$ inside $M_{\mathbb C}$ (the assumed \texttt{IsRealAnalytic},
  \texttt{HolomorphicComplexification}, \texttt{ComplexAnalyticStructure}). This is the
  paper's heaviest assumed piece. It also covers the totally-real $d$-submanifolds of
  $\mathbb M_{\mathbb C}$ used by the vacuum-domain candidate (folded in from the
  Minkowski geometry). The complexification of a real-analytic manifold is not in Mathlib.
  \emph{Research-scale.}
\end{definition}

\begin{definition}[The Wick / boundary correspondence]
  \label{found:wick-correspondence}
  \notready
  \uses{found:real-analytic-complexification}
  The boundary relation realizing $\mathcal C_d^{\mathrm{Lor},\omega}$ inside
  $\mathcal C_d^{\mathbb C}$: the object-level Wick rotation (a Lorentzian germ to the
  complex object it rotates to, the assumed \texttt{WickObjectCorrespondence}) and the
  morphism-level realization (a Lorentzian cobordism to a complex one, the assumed
  \texttt{CobordismRealization}). This is a construction KS defer, not a parametrization, so
  it is orange. \emph{Research-scale.}
\end{definition}

\begin{definition}[Hilbert completion of a pre-Hilbert space]
  \label{found:hilbert-completion}
  \lean{WickRotation.HilbertCompletion, WickRotation.hilbertCompletion}
  \leanok
  The completion of a positive-definite complex inner product space $E$ to a Hilbert space
  $E^{\mathrm{Hilb}}$, together with the dense injective embedding
  $\iota \colon E \hookrightarrow E^{\mathrm{Hilb}}$ and the inner-product extension
  $\langle \iota x, \iota y\rangle = \langle x, y\rangle$, built over Mathlib's completion
  machinery.
\end{definition}

\begin{definition}[The rigged triple and the $\kappa$-factorisation]
  \label{found:rigged-triple}
  \lean{WickRotation.InducesUnitaryGH.checkIncl, WickRotation.InducesUnitaryGH.hatIncl,
    WickRotation.InducesUnitaryGH.factor, WickRotation.InducesUnitaryGH.inner_eq,
    WickRotation.HilbertCompletion}
  \notready
  \uses{found:hilbert-completion, found:tvs-limits-tensor}
  The placement of the Hilbert space $E_\Sigma^{\mathrm{Hilb}}$ in the rigged triple
  $\check E_\Sigma \subset E_\Sigma^{\mathrm{Hilb}} \subset \hat E_\Sigma$: the downstream
  injective dense inclusion $E_\Sigma^{\mathrm{Hilb}} \subset \hat E_\Sigma$ and the factorisation
  of the composite $\check E_\Sigma \to \hat E_\Sigma$ through the field-theory map $\kappa$. The
  reflection-positivity pairing of Definition~\ref{def:unitarity} makes $\check E_\Sigma$ positive
  definite, so $E_\Sigma^{\mathrm{Hilb}}$ is the Hilbert completion
  (Definition~\ref{found:hilbert-completion}) of it. Here $\check E_\Sigma = \mathtt{T.ECheck}$ and
  $\hat E_\Sigma = \mathtt{T.EHat}$ are the deferred inverse and direct limit spaces
  (Definition~\ref{found:tvs-limits-tensor}), so there is no concrete $\hat E_\Sigma$ to map into;
  this wiring is the assumed \texttt{hatIncl}/\texttt{factor} data of \texttt{InducesUnitaryGH}.
\end{definition}

\begin{definition}[The observable bundle]
  \label{found:observable-bundle}
  \notready
  \uses{found:tvs-limits-tensor, def:metc-complex-manifold, found:smooth-cobordism-geometry}
  The observable bundle over $\mathrm{Met}_{\mathbb C}(\hat x)$ and the space of observables
  $\mathcal O_x = \varprojlim E_{\partial D}$ as the inverse limit over the boundary spheres
  of shrinking discs around $x$, with the disc-removal action
  $E_{\partial D} \to \mathrm{Hom}(E_{\Sigma_0}; E_{\Sigma_1})$ (the assumed
  \texttt{Observables.Obs}; the opaque ambient morphism has no disc-removal accessor).
  \emph{Research-scale.}
\end{definition}

\begin{definition}[Several complex variables: tube, Stein and Shilov theory]
  \label{found:scv-tube-domain}
  \notready
  Domains of holomorphy and holomorphic envelopes, Bochner's tube theorem, Siegel domains
  and the Cayley transform, the Shilov boundary, and the Wightman permuted extended tube
  $\mathcal U_k$ (the assumed \texttt{MinkowskiComplexGeometry.Uk} and
  \texttt{.isDomainOfHolomorphy}). This single body of several-complex-variable analysis is
  what Section~2 defers at Propositions~2.3, 2.7 and Lemma~2.8, what the two-copies
  Shilov-boundary statement needs, and what the closing $\mathcal V_k$ conjecture invokes.
  None of it is in Mathlib. \emph{Research-scale.}
\end{definition}

\begin{definition}[The Hodge star on exterior powers]
  \label{found:hodge-star}
  \lean{KontsevichSegal.Hodge.star_g, KontsevichSegal.Hodge.star_g_star_g}
  \leanok
  The Hodge star operator $*_g$ for complex metrics and the induced quadratic forms on the
  exterior powers $\Lambda^p(V^*)$, with the twisted determinant line. This is exactly the
  machinery Definition~2.1 (the Hodge-star form of allowability) and Theorem~2.2 (its
  equivalence with the angle condition) need, and the $p=0$ branch of the two-copies
  statement. The operator is built from scratch (Mathlib has no Hodge star at this
  generality): the perfect wedge pairing on exterior powers, the metric-normalized volume
  element, $*_g$ with $*_g*_g = (-1)^{p(d-p)}$, and the induced forms, are formalized in
  \texttt{ComplexMetrics/HodgeScaffold.lean}.
\end{definition}

\begin{definition}[Homogeneous spaces, associated bundles and the $Q_{\mathbb C}$ topology]
  \label{found:homogeneous-bundles}
  \notready
  The $\mathrm{GL}(V)/\mathrm O(V)$ and $\mathrm O_{\mathbb C}(Z)$ homogeneous-space actions,
  the associated-bundle machinery, the real Grassmannian, and a topology on
  $Q_{\mathbb C}(V)$ (needed even to state contractibility). This is what Propositions~2.4
  and~2.6 rest on, and what the two-dimensional polydisc description needs.
  \emph{Mathlib-scale.}
\end{definition}

\begin{lemma}[Courant--Fischer min--max and eigenvalue interlacing for allowable angles]
  \label{lem:eigenvalue-minmax}
  \notready
  \uses{def:allowable, def:restrict-form}
  For $g \in Q_{\mathbb C}(V)$, the ordered angles
  $\theta_1 \geq \dots \geq \theta_d$ of $g$ are characterized by the min--max
  \[
    \theta_k \;=\; \sup_{\dim A = k}\ \inf_{v \in \mathbb P(A)} \arg g(v)
             \;=\; \inf_{\dim A = d-k+1}\ \sup_{v \in \mathbb P(A)} \arg g(v).
  \]
  Consequently, if $W \leq V$ has codimension one, the angles
  $\theta'_1 \geq \dots \geq \theta'_{d-1}$ of $g|_W$ interleave those of $g$:
  $\theta_k \geq \theta'_k \geq \theta_{k+1}$.
\end{lemma}

\begin{definition}[Complexified Minkowski space]
  \label{found:minkowski-linear}
  \lean{WickRotation.MinkowskiLinear, WickRotation.MinkowskiModel.minkowskiLinear}
  \leanok
  Real Minkowski space $\mathbb M = \mathbb R^{d-1,1}$ with its Lorentzian form of signature
  $(d-1,1)$; its complexification
  $\mathbb M_{\mathbb C} = \mathbb M \otimes_{\mathbb R} \mathbb C = \mathbb M \oplus \mathrm i
  \mathbb M$, with the $\mathbb C$-bilinear base change of the real form; and the Euclidean real
  subspace $\mathbb E = \mathbb R^d$ of
  $\mathbb M_{\mathbb C} = \mathbb E \oplus \mathrm i \mathbb E$, the Wick rotation of $\mathbb M$,
  on which the $\mathbb C$-bilinear form is real and positive definite. This is the linear-algebra
  interface of \texttt{MinkowskiComplexGeometry} (its fields \texttt{mC}, \texttt{bilin},
  \texttt{Eucl}, \texttt{projE}), which now extends it. It is built as finite-dimensional linear
  algebra over Mathlib and realized by a concrete coordinate model on $\mathbb C^{d}$, $d \ge 2$:
  the restriction of the complex form to the real points is the real Lorentzian form, both
  splittings hold, the form is positive definite on $\mathbb E$, and the Lorentzian form is
  indefinite and nondegenerate; the forward light cone is exhibited (nonempty, open, proper), its
  convexity left to the causal geometry of
  Definition~\ref{found:lorentzian-causal-geometry}. The totally-real-submanifold half of
  \texttt{MinkowskiComplexGeometry} stays deferred, folded into
  Definition~\ref{found:real-analytic-complexification}.
\end{definition}
